\documentclass{article}
\usepackage[utf8]{inputenc}
\usepackage[T1]{fontenc}
\usepackage{graphicx}
\usepackage{algorithm}
\usepackage{algpseudocode}

\usepackage{listings}
\usepackage{xcolor}
\usepackage[margin=2cm]{geometry}
\usepackage{float}

\usepackage{amssymb} %Ajoute des symboles mathématiques comme R, Z, forall, exists, etc....
\usepackage{amsmath}
\usepackage{amsthm}


\lstdefinestyle{monstyle}{
    backgroundcolor=\color{white},
    commentstyle=\color{gray},
    keywordstyle=\color{blue},
    numberstyle=\tiny\color{gray},
    basicstyle=\ttfamily\small,
    breaklines=true,
    numbers=left,
    numbersep=5pt,
    frame=single,
}
\lstset{style=monstyle}


\title{\textbf{Rapport de TD n°6 : Initiation à l'algorithmique}}
\author{Derrez Lorenzo}

\begin{document}
\maketitle

\section*{\textbf{Exercice 0}}

\includegraphics[width=15cm]{exo0_TD6.png}

\section*{\textbf{Exercice 1}}

\subsection*{Question 1:}

\begin{itemize}
    \item Il faut des entiers long en entrée. Le cours définit le type Longint en Python I1 avec int64. 
    \item \textbf{Remarque :} Comme on modifie par effets de bords : Longint, on suppose que son allocation a été effectuée lors du main puis une fois que l'on retourne on libèrera la mémoire avec freeLongint.
\end{itemize}

\begin{lstlisting}[label=Python, caption={Algo1\_rec}]
def algo1_rec (n : Longint) -> Longint : 
    if n <= 10: 
        return 10
    else:
        if (n%2) != 0:
            return algo1_rec(n+3)+algo1_rec(n+1)
        else:
            return algo1_rec(n-2)+algo1_rec(n-7)+algo1_rec(shift(n,-1))
\end{lstlisting}

\subsection*{Question 2}

\begin{itemize}
    \item Evaluation de la complexité en espace : 
    \subitem Chacune des variables occupe une place constante 
    \subitem Il y a des appels récursifs $\implies$ augmente la place occupée sur la pile d'appels $\implies$ augmente la complexité en espace.
    \subitem Globalement (avec le shift), on descend de 2 en 2 $\implies \frac{n}{2}$ appels empilés $\implies \#espace = \theta(n)$
    \item Evaluation de la complexité en temps : 
    \subitem Au cours de chaque appel, on peut faire :
        \subsubitem une comparaison 
        \subsubitem un modulo 
        \subsubitem des affectations 
        \subsubitem des appels récursifs 
\end{itemize}

\textbf{Pour la complexité en temps : } (arbre d'appels) 

\includegraphics[width=15cm]{img_test_Q2_TD6.png} etc...

\begin{itemize}
    \item Chaque appel engendre plusieurs sous calculs sur des arguments qui diminuent lentement. (n-2, n-4, etc...)
    \item Cela entraîne une \textbf{complexité en temps exponentielle.}
\end{itemize}

\subsection*{Question 3}

\begin{itemize}
    \item une des solutions pour éviter de calculer une seconde fois une valeur est la \textbf{mémoisation}
    \item memo : tableau de (n+4) initialiée à -1.
\end{itemize}

\begin{lstlisting}[label=Python, caption={Algo2\_rec}]
def algo2_rec(memo : Array[integer], n : Longint) -> Longint : 
    if memo[n] != -1:
        return memo[n]
    if n<=10:
        resultat = 10
    elif (n%2) != 0:
        resultat = algo2_rec(memo, n+1)+algo2_rec(memo, n+3)
    else: 
        resultat = algo2_rec(memo, n-2)+algo2_rec(memo, n-7)+algo2_rec(memo, shift(n, -1))
    memo[n] = resultat
    return resultat
\end{lstlisting}

\subsection*{Question 4}

\includegraphics[width=15cm]{image_Q4_exo1_TD6.png}

\begin{itemize}
    \item Complexité en espace : 
    \subitem $ \#f(n) = \theta(n)_{pile memoire} + \theta(n\log n)_{memo} = \theta(n) $
    \item Complexité en temps : 
    \subitem $ \theta(n) $ malgré les appels récursifs car memo permet de récupérer directemen les valeurs déjà calculées.
\end{itemize}

\section*{\textbf{Exercice 2}}

\subsection*{Question 1}

\begin{itemize}
    \item Le problème est syntaxiquement proche du problème \textit{LongeurMaxPlateau}.
\end{itemize}

\subsection*{Question 2}

\begin{itemize}
    \item Principe : 
    \subitem tri en $\theta(n\log n) < \theta(n^2) $ 
    \subitem parcours unique en $\theta(n)$ 
\end{itemize}

\begin{lstlisting}[label=python, caption={isDiversity}]
def isDiversity(i : integer, j : integer, T : Array[integer]) -> bool : 
    S = Array.copy(T[i,j])
    trifusion(S)

    for k in range(1, len(S)):
        if S[k] ==  S[k-1]:
            return false
    return true
\end{lstlisting}

\subsection*{Question 3}

\begin{itemize}
    \item Evaluation de la complexité temporelle : 
    \subitem $ \# f(n) = \theta(n)_{copy} + \theta(n \log n)_{trifusion} + \theta(n)_{parcours unique} = \theta(n \log n) < \theta(n^2) $
\end{itemize}

\subsection*{Question 4}

\begin{lstlisting}[label=python, caption={max\_variety}]
def max_variety(T: Array[integer]) -> int:
    curr_max = 1
    max = 0
    curr_i = 0
    curr_j = 1

    while curr_i != len(T):
        different = False

        # Verifier que le nouvel element est different de tous les elements du sous-tableau courant.
        if curr_j < len(T):
            different = True
            for k in range(curr_i, curr_j):
                if T[k] == T[curr_j]:
                    different = False
                    break

        if different:
            curr_max = curr_max + 1
            curr_j = curr_j + 1
        else:
            if curr_max > max:
                max = curr_max

            curr_i = curr_i + 1
            curr_j = curr_i + 1
            curr_max = 1

    return max
\end{lstlisting}

\subsection*{Question 5}

\begin{itemize}
    \item Evaluation de la complexité en espace : 
    \subitem affectation, introduction de variables entières, comparaisons $ \implies \theta(1) $
    \item Evaluation de la complexité en temps : 
    \subitem Parcours de curr\_i à len(T) $ \implies \theta(n) $ (où n = len(T) - curr\_i) 
    \subitem Au sein de la boucle dans le pire des cas, on vérifie tous les cas possibles pour chaque indice $\theta(n^2)$ 
    \subitem $ \implies \theta(n^3) $ au total. 
\end{itemize}

\subsection*{Question 6}

\begin{lstlisting}[label=python, caption={\textit{nextIndex}}]
def nextIndex(T):
    n = len(T)
    U = [n] * n

    i = 0
    while i < n:
        for j in range(n - 1, i, -1):
            if T[i] == T[j]:
                U[i] = j
        i += 1

    return U
\end{lstlisting}

\begin{itemize}
    \item \textbf{Remarque :} Cette méthode a une complexité temporelle $\theta(n^2)$ et en espace $\theta(n)$. 
    \subitem Elle n'est peut être pas optimale. Une solution à l'aide d'un dictionnaire pourrait potentiellement éviter les doubles boucles imbriquées. Dès lors on aurait une complexité en $\theta(n)$ 
\end{itemize}

\subsection*{Question 7}

On considère le tableau de longueur $n=20$ suivant :
\[
\begin{aligned}
T = [&1,2,1,3,4,2,5,6,3,7,4,8,5,2,6,1,7,3,8,4].
\end{aligned}
\]

Son tableau prochain indice est :
\[
\begin{aligned}
U = [&2,5,15,8,10,13,12,14,17,16,19,18,20,20,20,20,20,20,20,20].
\end{aligned}
\]
Les indices commencent à $0$. La valeur $20$ indique
qu'il n'existe aucune occurrence suivante.
Par exemple, $U[3]=8$ car $T[3]=T[8]=3$,
sans autre occurrence de $3$ entre ces deux indices.

Le sous-tableau
\[
T[8:16]=[3,7,4,8,5,2,6,1]
\]
contient huit éléments deux à deux différents :
c'est donc une diversité de longueur $8$.

De plus, $T$ ne contient que huit valeurs distinctes,
les entiers de $1$ à $8$. Aucune diversité ne peut
donc avoir une longueur supérieure à $8$.

Ainsi, la sortie de la fonction est :
\[
\boxed{\operatorname{variety}(T)=8}.
\]

\subsection*{Question 8}



\end{document}